\documentclass[12pt,a4paper]{article}
\usepackage{../../../myconfig}

\begin{document}

% =========================================================================
% BÀI TẬP ĐIỂM DANH: XÁC SUẤT CƠ BẢN VÀ XÁC SUẤT CỔ ĐIỂN (10 CÂU)
% Dạng trả lời ngắn. Đáp án được ghi dưới từng câu bằng comment.
% =========================================================================
\titlebox{Xác suất}{Xác suất cơ bản \\ Xác suất cổ điển}{Xác suất: Xác suất cơ bản và xác suất cổ điển}

\vspace{-0.25cm}

% --- Câu 1 (Câu 6 trong tài liệu gốc) ---
\questionShortAnswer{0}{1}{%
    Gieo một con xúc xắc cân đối và đồng chất một lần.
    Gọi $A$ là biến cố ``xuất hiện mặt 1 chấm'', $B$ là biến cố ``xuất hiện mặt có số chấm lẻ''.
    Tính $P(A\mid B)$ (làm tròn đến hàng phần trăm).
}
\drawdashlinesfull
% Đáp án: 0,33.
% Kiểm tra: P(A|B) = (1/6)/(3/6) = 1/3.

% --- Câu 2 (Câu 7 trong tài liệu gốc) ---
\questionShortAnswer{0}{2}{%
    Một phòng thí nghiệm thực hiện hai thí nghiệm độc lập.
    Xác suất thành công của mỗi thí nghiệm là $0{,}8$.
    Phòng được coi là hoàn thành nhiệm vụ nếu có ít nhất một thí nghiệm thành công.
    Tính xác suất phòng hoàn thành nhiệm vụ (làm tròn đến hàng phần trăm).
}
\drawdashlinesfull
% Đáp án: 0,96.
% Kiểm tra: P = 1 - (1 - 0,8)^2 = 0,96.


% --- Câu 3 (Câu 8 trong tài liệu gốc) ---
\questionShortAnswer{0}{3}{%
    Một máy ép thủy lực có hai động cơ $A$ và $B$ hoạt động độc lập.
    Xác suất động cơ $A$ chạy tốt là $0{,}8$; xác suất động cơ $B$ chạy tốt là $0{,}75$.
    Máy hoạt động nếu có ít nhất một động cơ chạy tốt.
    Tính xác suất máy hoạt động (làm tròn đến hàng phần trăm).
}
\drawdashlinesfull
% Đáp án: 0,95.
% Kiểm tra: P = 1 - (1 - 0,8)(1 - 0,75) = 0,95.


% --- Câu 4 (Câu 11 trong tài liệu gốc) ---
\questionShortAnswer{0}{4}{%
    Một hộp đựng 7 quả cầu trắng và 3 quả cầu đỏ, các quả cầu có cùng kích thước và khối lượng.
    Lấy ngẫu nhiên đồng thời 4 quả cầu.
    Tính xác suất trong 4 quả lấy được có đúng 2 quả đỏ (làm tròn đến hàng phần trăm).
}
\drawdashlinesfull
% Đáp án: 0,30.
% Kiểm tra: P = C_3^2 C_7^2 / C_{10}^4 = 3/10.


% --- Câu 5 (Câu 12 trong tài liệu gốc) ---
\questionShortAnswer{0}{5}{%
    Một bình có 50 viên bi cùng kích thước và chất liệu, gồm 30 bi đen và 20 bi trắng.
    Lấy ngẫu nhiên lần lượt 2 viên bi, không hoàn lại.
    Tính xác suất lấy được bi đen ở lần thứ nhất và bi trắng ở lần thứ hai (làm tròn đến hàng phần trăm).
}
\drawdashlinesfull
% Đáp án: 0,24.
% Kiểm tra: P = (30/50)(20/49) = 12/49.


% --- Câu 6 (Câu 4 trong tài liệu gốc) ---
\questionShortAnswer{0}{6}{%
    Một công ty đấu thầu hai dự án.
    Xác suất thắng dự án 1 và dự án 2 lần lượt là $50\%$ và $60\%$; xác suất thắng cả hai là $40\%$.
    Biết công ty không thắng dự án 1, tính xác suất công ty thắng dự án 2 (làm tròn đến hàng phần trăm).
}
\drawdashlinesfull
% Đáp án: 0,40.
% Kiểm tra: P = (0,6 - 0,4)/(1 - 0,5) = 0,4.


% --- Câu 7 (Câu 1 trong tài liệu gốc) ---
\questionShortAnswer{0}{7}{%
    Hai người gọi điện thoại đến hai số khác nhau nhưng đều quên chữ số cuối.
    Mỗi người thử ngẫu nhiên các chữ số từ 0 đến 9, không lặp lại chữ số đã thử; hai người thử độc lập.
    Tính xác suất ít nhất một người gọi đúng mà không phải thử quá hai lần (làm tròn đến hàng phần trăm).
}
\drawdashlinesfull
% Đáp án: 0,36.
% Kiểm tra: P = 1 - (1 - 2/10)^2 = 0,36.


% --- Câu 8 (Câu 3 trong tài liệu gốc) ---
\questionShortAnswer{0}{8}{%
    Hộp I có 6 bi đỏ và 4 bi xanh; hộp II có 5 bi đỏ và 5 bi xanh.
    Các viên bi có cùng kích thước và khối lượng.
    Minh lấy ngẫu nhiên một bi từ hộp I; Như lấy ngẫu nhiên một bi từ hộp II, độc lập với Minh.
    Biết có ít nhất một bi đỏ được lấy ra, xác suất Như lấy được bi đỏ là $\dfrac ab$ (phân số tối giản, $a,b$ nguyên dương).
    Tính $a-b$.
}
\drawdashlinesfull
% Đáp án: -3.
% Kiểm tra: P = (1/2)/(1 - (4/10)(5/10)) = 5/8; a - b = 5 - 8 = -3.


% --- Câu 9 (Câu 17 trong tài liệu gốc) ---
\questionShortAnswer{0}{9}{%
    Một người môi giới có 8 chìa khóa ứng với 8 ngôi nhà, mỗi chìa chỉ mở được đúng một nhà.
    Có 3 nhà không khóa cửa.
    Người đó chọn ngẫu nhiên 3 chìa mang theo, rồi chọn ngẫu nhiên một trong 8 nhà để vào, độc lập với việc chọn chìa.
    Tính xác suất người đó vào được nhà (làm tròn đến hàng phần trăm).
}
\drawdashlinesfull
% Đáp án: 0,61.
% Kiểm tra: P = 3/8 + (5/8)(3/8) = 39/64.


% --- Câu 10 (Câu 18 trong tài liệu gốc) ---
\questionShortAnswer{0}{10}{%
    Một trường có 8 giáo viên Toán gồm 3 nữ và 5 nam, cùng 4 giáo viên Vật lý đều là nam.
    Chọn ngẫu nhiên một đoàn 3 người.
    Tính xác suất đoàn có đủ hai môn Toán và Vật lý, đồng thời có cả nam và nữ (làm tròn đến hàng phần trăm).
}
\drawdashlinesfull
% Đáp án: 0,41.
% Kiểm tra: P = [C_4^2 C_3^1 + C_4^1(C_8^2 - C_5^2)]/C_{12}^3 = 9/22.

\end{document}
